% Historical ADMM derivation retained as a style reference.
% Source: archival/supplement.tex at commit b2c83675f3a1cc78dcc4696a9aa85279927c4288.
% This file is not included in the manuscript.
% Preserve this historical copy; revise the active proof in supplement.tex.

\subsubsection{Proof of Proposition \ref{prop:computational_complexity}}
\label{proof:computational_complexity}

\begin{restatedproposition}{prop:computational_complexity}
The computational complexity of each step of Isometry Pursuit is $O(D^2P)$.
\end{restatedproposition}

% TODO: Replace the direct-ADMM derivation below with the SCS cached linear-solve and cone-projection analysis for the runtime panel of Figure 3, distinguishing setup from per-iteration work.
\begin{proof}
The augmented split Lagrangian minimized by ADMM is
    \begin{align}
    \mathcal L_{\rho,\gamma, \nu} (X,\beta,z) = \|z\|_{1,2} + \text{tr}(\nu^T (I_D - n_q^{\times P}(X) \beta)) + \rho \|I_D - n_q^{\times P}(X) \beta\|_F^2 + \gamma \|\beta - z\|_F^2
    \end{align}
This minimization problem is
    \begin{align}
        \arg \min_{\beta,z \in P \times D} \mathcal L_{\rho,\gamma,\nu} (X,\beta,z)
    \end{align}
For a more detailed explanation of the reasons to prefer optimizing the split augmented Lagrangian over direct optimization of Equation \ref{prog:isometry_pursuit} see \citet{Boyd2010-jd}.
    
The three main subroutines ( \ref{eq:beta_update} - \ref{eq:nu_update}) called by ADMM determine the per iteration complexity.
The update rules of the ADMM algorithm used to optimize the augmented split Lagrangian are
    \begin{align}
    \beta^{k+1} = \arg \min_{\beta \in P \times D} \mathcal L_{\rho,\gamma,\nu} (X,\beta,z^k) \\
    z^{k+1} = \arg \min_{z \in P \times D} \mathcal L_{\rho,\gamma,\nu} (X,\beta^{k+1},z) \\
    \nu^{k+1} = \nu^k + \rho (I_D - n_q^{\times P}(X) \beta^{k+1})
    \end{align}
We will analyze the computational complexity of each of these steps by rewriting the update rules according to the specifics of the isometry pursuit problem.

The beta update Equation \ref{eq:beta_update} projects the solution onto the feasible set.
\begin{align}
    \label{eq:beta_update}
    \beta^{k+1} &= \arg \min_{\beta \in P \times D}  \text{tr}({\nu^k}^T (I_D - n_q^{\times P}(X) \beta)) + \rho \|I_D - n_q^{\times P}(X) \beta\|_2^2 + \gamma \|\beta - z^k\|_2^2 \\
    &= (\rho \, n_q^{\times P}(X)^T n_q^{\times P}(X) + \gamma I_P)^{-1} (\rho \, n_q^{\times P}(X)^T I_D + \gamma z^k + n_q^{\times P}(X)^T \nu^k)
    \end{align}
% This is accomplished through computation of the pseudo-inverse of the linear regression solution $(n_q^{\times P}(X)^T n_q^{\times P}(X))^{-1} n_q^{\times P}(X)^T I_D$.
Write $N=n_q^{\times P}(X)$.
The Woodbury identity gives
\begin{align}
(\gamma I_P+\rho N^\top N)^{-1}
=\gamma^{-1}\left[I_P-\rho N^\top
(\gamma I_D+\rho NN^\top)^{-1}N\right].
\end{align}
For fixed $N$, $\rho>0$, and $\gamma>0$, forming $NN^\top$ and factoring $\gamma I_D+\rho NN^\top$ costs $O(PD^2+D^3)$ once.
To apply the inverse to a $P\times D$ right-hand side $R$, compute $NR$, solve $(\gamma I_D+\rho NN^\top)Z=NR$ using the cached factorization, and return $\gamma^{-1}(R-\rho N^\top Z)$.
The two matrix products cost $O(PD^2)$ and the solve costs $O(D^3)$.
Forming the right-hand side of the coefficient update also costs $O(PD^2)$, so the per-iteration cost is $O(PD^2)$ when $P\geq D$.

Subsequently, the z update Equation \ref{eq:z_update} is soft-thresholding operation onto a contracted level set $\|\|_{1,2}$ ball.
\begin{align}
    \label{eq:z_update}
    z^{k+1} &= \arg \min_{z \in P \times D} \|z\|_{1,2} + \gamma \|z - \beta^{k+1}\|_2^2 \\
    &= \begin{cases} (1 - \frac{1}{\gamma\|\beta_{j.}^{k+1}\|_2}) (\beta_{j.}^{k+1}) \; \text{if} \; \|\beta_{j.}^{k+1}\|_2 > \frac{1}{\gamma}\\
        0 \; \text{otherwise}
    \end{cases}
\end{align}
The computational complexity of this operation is $O(PD)$.

Finally, the $\nu$ update Equation \ref{eq:nu_update} increases the squeezing of the $z$ and $\beta$ variables towards to each other.
\begin{align}
\nu^{k+1} = \nu^k + \rho (I_D - n_q^{\times P}(X) \beta^{k+1})
\label{eq:nu_update}
\end{align}
In this step, the dual variable $\nu$ that controls the penalty accrued by the splitting is updated in the direction of the constraint violation.
The computational complexity of this operation is dominated by that of the matrix multiplication $n_q^{\times P}(X) \beta^{k+1}$.
That is, it is $O(D^2P)$.
\end{proof}

